\documentclass[11pt,a4paper]{article}

\usepackage{iftex}
\ifPDFTeX
    \IfFileExists{lato.sty}{
        \usepackage[default]{lato}
    }{
        \usepackage{helvet}
        \renewcommand{\familydefault}{\sfdefault}
    }
    \usepackage[T1]{fontenc}
    \usepackage[utf8]{inputenc}
\else
    \usepackage{fontspec}
    \IfFontExistsTF{Lato}{
        \setmainfont{Lato}
        \setsansfont{Lato}
    }{
        \setmainfont{TeX Gyre Heros}
        \setsansfont{TeX Gyre Heros}
    }
\fi

\usepackage[a4paper,margin=1in]{geometry}
\usepackage{amsmath,amssymb,mathtools}
\usepackage{booktabs}
\usepackage{enumitem}
\usepackage{microtype}
\usepackage{xcolor}
\usepackage{titlesec}
\usepackage{hyperref}

\definecolor{accent}{HTML}{1F6F8B}
\definecolor{textgray}{HTML}{333333}

\hypersetup{
    colorlinks=true,
    linkcolor=accent,
    citecolor=accent,
    urlcolor=accent
}

\titleformat{\section}
    {\Large\bfseries\color{accent}}
    {\thesection}{0.75em}{}
\titleformat{\subsection}
    {\large\bfseries\color{textgray}}
    {\thesubsection}{0.75em}{}

\setlist[itemize]{leftmargin=1.6em,itemsep=0.2em}
\setlist[enumerate]{leftmargin=1.6em,itemsep=0.2em}

\title{\textbf{A Bounded-Trip MILP Model for the MPVRP-CC}}
\author{MPVRP-CC Computational Model}
\date{May 2026}

\begin{document}

\maketitle

\begin{abstract}
This document states the mixed-integer linear programming model implemented in
\texttt{mpvrp\_cc/optimization/milp\_solver.py} for the Multi-Product Vehicle Routing Problem with Split
Deliveries and Changeover Costs (MPVRP-CC). Each vehicle may perform a bounded
number of mini-routes. A mini-route starts at one depot, loads one product,
visits one or more stations, and ends at one depot. The model minimizes travel
distance, empty depot transfers, garage access and return distance, and product
changeover costs. The formulation also enforces depot stock limits and prevents
the same vehicle from visiting the same station more than once for the same
product.
\end{abstract}

\section{Problem Setting}

Let \(R\) be the maximum number of mini-routes allowed for each vehicle. In the
implementation, this value is the argument
\texttt{max\_trips\_per\_vehicle}. If it is not provided, the code computes the
lower bound
\[
    R =
    \left\lceil
    \frac{\sum_{s \in S}\sum_{p \in P} H_{sp}}
         {\sum_{k \in K} Q_k}
    \right\rceil ,
\]
where \(H_{sp}\) is the demand of station \(s\) for product \(p\), and \(Q_k\)
is the capacity of vehicle \(k\).

The model allows split deliveries. Therefore, the demand of a station for a
product may be served by several vehicles or by several trips, subject to the
additional rule that one vehicle cannot visit the same station twice for the
same product.

\section{Sets and Indices}

\begin{table}[h]
\centering
\begin{tabular}{p{0.22\linewidth}p{0.68\linewidth}}
\toprule
Symbol & Description \\
\midrule
\(K\) & Set of vehicles, indexed by \(k\). \\
\(T=\{1,\ldots,R\}\) & Ordered set of mini-route positions, indexed by \(t\). \\
\(P\) & Set of products, indexed by \(p\). \\
\(D\) & Set of depots, indexed by \(d\). \\
\(S\) & Set of stations, indexed by \(s\). \\
\(N=D \cup S\) & Set of nodes used inside a mini-route. \\
\(A\) & Set of feasible directed arcs \((i,j)\), \(i,j \in N\), \(i\neq j\), excluding depot-to-depot arcs. \\
\bottomrule
\end{tabular}
\end{table}

Depot-to-depot movements are not internal mini-route arcs. They are represented
separately between two consecutive mini-routes of the same vehicle.

\section{Parameters}

\begin{table}[h]
\centering
\begin{tabular}{p{0.22\linewidth}p{0.68\linewidth}}
\toprule
Symbol & Description \\
\midrule
\(Q_k\) & Capacity of vehicle \(k\). \\
\(H_{sp}\) & Demand of station \(s\) for product \(p\). \\
\(B_{dp}\) & Stock of product \(p\) available at depot \(d\). \\
\(g(k)\) & Home garage of vehicle \(k\). \\
\(p^0_k\) & Initial product configuration of vehicle \(k\). \\
\(C_{pp'}\) & Changeover cost from product \(p\) to product \(p'\). \\
\(\Delta_{ij}\) & Distance between nodes \(i\) and \(j\). \\
\(M=|S|\) & Big-\(M\) value used for station ordering. \\
\bottomrule
\end{tabular}
\end{table}

\section{Decision Variables}

\begin{table}[h]
\centering
\begin{tabular}{p{0.18\linewidth}p{0.18\linewidth}p{0.54\linewidth}}
\toprule
Variable & Domain & Meaning \\
\midrule
\(a_{kt}\) & \(\{0,1\}\) & 1 if vehicle \(k\) uses trip \(t\). \\
\(y_{ktp}\) & \(\{0,1\}\) & 1 if trip \((k,t)\) carries product \(p\). \\
\(b_{ktd}\) & \(\{0,1\}\) & 1 if trip \((k,t)\) starts at depot \(d\). \\
\(e_{ktd}\) & \(\{0,1\}\) & 1 if trip \((k,t)\) ends at depot \(d\). \\
\(v_{kts}\) & \(\{0,1\}\) & 1 if trip \((k,t)\) visits station \(s\). \\
\(w_{ktsp}\) & \(\{0,1\}\) & 1 if trip \((k,t)\) visits station \(s\) while carrying product \(p\). \\
\(x_{ktij}\) & \(\{0,1\}\) & 1 if arc \((i,j)\in A\) is used in trip \((k,t)\). \\
\(q_{ktsp}\) & \(\mathbb{R}_+\) & Quantity of product \(p\) delivered to station \(s\). \\
\(L_{ktdp}\) & \(\mathbb{R}_+\) & Quantity of product \(p\) loaded from depot \(d\). \\
\(u_{kts}\) & \([0,M]\) & Order of station \(s\) in trip \((k,t)\). \\
\(m_{kt}\) & \(\{0,1\}\) & 1 if trip \(t\) is the last active trip of vehicle \(k\). \\
\(r_{ktd}\) & \(\{0,1\}\) & 1 if vehicle \(k\) returns from depot \(d\) after trip \(t\). \\
\(z_{ktpp'}\) & \(\{0,1\}\) & 1 if vehicle \(k\) changes from product \(p\) to \(p'\) after trip \(t\). \\
\(f_{ktdd'}\) & \(\{0,1\}\) & 1 if vehicle \(k\) transfers from depot \(d\) to depot \(d'\) between trips. \\
\bottomrule
\end{tabular}
\end{table}

\section{Objective Function}

The objective minimizes the total operating cost:
\begin{align}
\min \quad
& \sum_{k \in K}\sum_{t \in T}\sum_{(i,j)\in A}
    \Delta_{ij} x_{ktij}
    && \text{mini-route travel} \nonumber\\
& + \sum_{k \in K}\sum_{d \in D}
    \Delta_{g(k)d} b_{k1d}
    && \text{travel from garage to first depot} \nonumber\\
& + \sum_{k \in K}\sum_{t \in T}\sum_{d \in D}
    \Delta_{dg(k)} r_{ktd}
    && \text{return from last depot to garage} \nonumber\\
& + \sum_{k \in K}\sum_{t=1}^{R-1}\sum_{d \in D}\sum_{d' \in D}
    \Delta_{dd'} f_{ktdd'}
    && \text{empty transfer between depots} \nonumber\\
& + \sum_{k \in K}\sum_{p \in P}
    C_{p^0_k p} y_{k1p}
    && \text{initial product setup} \nonumber\\
& + \sum_{k \in K}\sum_{t=1}^{R-1}\sum_{p \in P}\sum_{p' \in P}
    C_{pp'} z_{ktpp'}
    && \text{inter-trip changeover}.
\label{eq:objective}
\end{align}

\section{Constraints}

\subsection{Trip Activation and Sequencing}

Each active trip selects exactly one product, one start depot, and one end
depot:
\begin{align}
    \sum_{p \in P} y_{ktp} &= a_{kt}
        && \forall k \in K,\ t \in T, \label{eq:one-product}\\
    \sum_{d \in D} b_{ktd} &= a_{kt}
        && \forall k \in K,\ t \in T, \label{eq:one-start}\\
    \sum_{d \in D} e_{ktd} &= a_{kt}
        && \forall k \in K,\ t \in T. \label{eq:one-end}
\end{align}

Trips are used in order:
\begin{equation}
    a_{kt} \leq a_{k,t-1}
    \qquad \forall k \in K,\ t=2,\ldots,R.
\label{eq:ordered-trips}
\end{equation}

\subsection{Mini-Route Flow}

If trip \((k,t)\) starts at depot \(d\), exactly one arc leaves that depot
toward a station:
\begin{equation}
    \sum_{j \in S} x_{ktdj} = b_{ktd}
    \qquad \forall k \in K,\ t \in T,\ d \in D.
\label{eq:start-flow}
\end{equation}

If trip \((k,t)\) ends at depot \(d\), exactly one station-to-depot arc enters
that depot:
\begin{equation}
    \sum_{i \in S} x_{ktid} = e_{ktd}
    \qquad \forall k \in K,\ t \in T,\ d \in D.
\label{eq:end-flow}
\end{equation}

For each station, incoming and outgoing flow are both equal to the visit
decision:
\begin{align}
    \sum_{\substack{i \in N\\i\neq s}} x_{ktis} &= v_{kts}
        && \forall k \in K,\ t \in T,\ s \in S, \label{eq:station-in}\\
    \sum_{\substack{j \in N\\j\neq s}} x_{ktsj} &= v_{kts}
        && \forall k \in K,\ t \in T,\ s \in S. \label{eq:station-out}
\end{align}

Let \(P_s=\{p\in P:H_{sp}>0\}\) be the set of products demanded by station
\(s\). A station can be visited only by a trip carrying a demanded product:
\begin{equation}
    v_{kts} \leq \sum_{p \in P_s} y_{ktp}
    \qquad \forall k \in K,\ t \in T,\ s \in S: P_s \neq \emptyset.
\label{eq:visit-product}
\end{equation}
If \(P_s=\emptyset\), then
\begin{equation}
    v_{kts}=0
    \qquad \forall k \in K,\ t \in T.
\label{eq:no-demand-visit}
\end{equation}

\subsection{Station-Product Visit Logic}

The variable \(w_{ktsp}\) is the logical AND of the station visit variable
\(v_{kts}\) and the selected product variable \(y_{ktp}\):
\begin{align}
    w_{ktsp} &\leq v_{kts}
        && \forall k \in K,\ t \in T,\ s \in S,\ p \in P, \label{eq:w-visit}\\
    w_{ktsp} &\leq y_{ktp}
        && \forall k \in K,\ t \in T,\ s \in S,\ p \in P, \label{eq:w-product}\\
    w_{ktsp} &\geq v_{kts} + y_{ktp} - 1
        && \forall k \in K,\ t \in T,\ s \in S,\ p \in P. \label{eq:w-and}
\end{align}

The same vehicle cannot visit the same station more than once for the same
product:
\begin{equation}
    \sum_{t \in T} w_{ktsp} \leq 1
    \qquad \forall k \in K,\ s \in S,\ p \in P.
\label{eq:one-visit-station-product}
\end{equation}

\subsection{Demand Satisfaction and Capacity}

All positive station-product demands must be satisfied:
\begin{equation}
    \sum_{k \in K}\sum_{t \in T} q_{ktsp} = H_{sp}
    \qquad \forall s \in S,\ p \in P: H_{sp}>0.
\label{eq:demand}
\end{equation}

No quantity is delivered for a zero-demand station-product pair:
\begin{equation}
    q_{ktsp}=0
    \qquad \forall k \in K,\ t \in T,\ s \in S,\ p \in P: H_{sp}=0.
\label{eq:zero-demand}
\end{equation}

Delivery quantities are linked to the selected product and station visit:
\begin{align}
    q_{ktsp} &\leq H_{sp} y_{ktp}
        && \forall k \in K,\ t \in T,\ s \in S,\ p \in P, \label{eq:quantity-product}\\
    q_{ktsp} &\leq H_{sp} v_{kts}
        && \forall k \in K,\ t \in T,\ s \in S,\ p \in P. \label{eq:quantity-visit}
\end{align}

Vehicle capacity is enforced for each trip:
\begin{equation}
    \sum_{s \in S}\sum_{p \in P} q_{ktsp}
    \leq Q_k a_{kt}
    \qquad \forall k \in K,\ t \in T.
\label{eq:capacity}
\end{equation}

\subsection{Depot Stock Accounting}

The model uses \(L_{ktdp}\) to assign the load of each trip to the selected
start depot. For each vehicle \(k\), let the trip load of product \(p\) be
\[
    Q^{\mathrm{trip}}_{ktp} = \sum_{s \in S} q_{ktsp}.
\]
The following linear constraints enforce
\(L_{ktdp}=Q^{\mathrm{trip}}_{ktp}\) if trip \((k,t)\) starts at depot \(d\),
and \(L_{ktdp}=0\) otherwise:
\begin{align}
    L_{ktdp} &\leq Q_k b_{ktd}
        && \forall k \in K,\ t \in T,\ d \in D,\ p \in P, \label{eq:load-start}\\
    L_{ktdp} &\leq \sum_{s \in S} q_{ktsp}
        && \forall k \in K,\ t \in T,\ d \in D,\ p \in P, \label{eq:load-upper}\\
    L_{ktdp} &\geq \sum_{s \in S} q_{ktsp} - Q_k(1-b_{ktd})
        && \forall k \in K,\ t \in T,\ d \in D,\ p \in P. \label{eq:load-lower}
\end{align}

The total quantity loaded from each depot cannot exceed the available stock:
\begin{equation}
    \sum_{k \in K}\sum_{t \in T} L_{ktdp}
    \leq B_{dp}
    \qquad \forall d \in D,\ p \in P.
\label{eq:stock}
\end{equation}

\subsection{Subtour Elimination}

The implementation uses Miller--Tucker--Zemlin station order variables:
\begin{align}
    u_{kts} &\leq M v_{kts}
        && \forall k \in K,\ t \in T,\ s \in S, \label{eq:order-upper}\\
    u_{kts} &\geq v_{kts}
        && \forall k \in K,\ t \in T,\ s \in S. \label{eq:order-lower}
\end{align}

For station-to-station arcs, subtours are removed by:
\begin{equation}
    u_{kti} - u_{ktj} + M x_{ktij} \leq M - 1
    \qquad
    \forall k \in K,\ t \in T,\ i,j \in S,\ i\neq j.
\label{eq:mtz}
\end{equation}

\subsection{Changeover and Depot Transfer Logic}

For consecutive trips, \(z_{ktpp'}\) is the logical AND of \(y_{ktp}\) and
\(y_{k,t+1,p'}\):
\begin{align}
    z_{ktpp'} &\leq y_{ktp}
        && \forall k \in K,\ t=1,\ldots,R-1,\ p,p' \in P, \label{eq:change-prev}\\
    z_{ktpp'} &\leq y_{k,t+1,p'}
        && \forall k \in K,\ t=1,\ldots,R-1,\ p,p' \in P, \label{eq:change-next}\\
    z_{ktpp'} &\geq y_{ktp}+y_{k,t+1,p'}-1
        && \forall k \in K,\ t=1,\ldots,R-1,\ p,p' \in P. \label{eq:change-and}
\end{align}

Similarly, \(f_{ktdd'}\) is the logical AND of ending trip \(t\) at depot \(d\)
and starting trip \(t+1\) at depot \(d'\):
\begin{align}
    f_{ktdd'} &\leq e_{ktd}
        && \forall k \in K,\ t=1,\ldots,R-1,\ d,d' \in D, \label{eq:transfer-end}\\
    f_{ktdd'} &\leq b_{k,t+1,d'}
        && \forall k \in K,\ t=1,\ldots,R-1,\ d,d' \in D, \label{eq:transfer-start}\\
    f_{ktdd'} &\geq e_{ktd}+b_{k,t+1,d'}-1
        && \forall k \in K,\ t=1,\ldots,R-1,\ d,d' \in D. \label{eq:transfer-and}
\end{align}

\subsection{Last Active Trip and Garage Return}

The last active trip of each vehicle is identified by:
\begin{align}
    m_{kt} &= a_{kt}-a_{k,t+1}
        && \forall k \in K,\ t=1,\ldots,R-1, \label{eq:last-trip-middle}\\
    m_{kR} &= a_{kR}
        && \forall k \in K. \label{eq:last-trip-final}
\end{align}

The variable \(r_{ktd}\) is the logical AND of ending trip \((k,t)\) at depot
\(d\) and trip \(t\) being the last active trip:
\begin{align}
    r_{ktd} &\leq e_{ktd}
        && \forall k \in K,\ t \in T,\ d \in D, \label{eq:return-end}\\
    r_{ktd} &\leq m_{kt}
        && \forall k \in K,\ t \in T,\ d \in D, \label{eq:return-last}\\
    r_{ktd} &\geq e_{ktd}+m_{kt}-1
        && \forall k \in K,\ t \in T,\ d \in D. \label{eq:return-and}
\end{align}

\section{Domains}

\begin{align}
    a_{kt}, y_{ktp}, b_{ktd}, e_{ktd}, v_{kts}, w_{ktsp}, x_{ktij},
    m_{kt}, r_{ktd}, z_{ktpp'}, f_{ktdd'} &\in \{0,1\}, \\
    q_{ktsp}, L_{ktdp} &\geq 0, \\
    0 \leq u_{kts} &\leq M.
\end{align}

\section{Implementation Notes}

The notation above follows the implemented model in
\texttt{mpvrp\_cc/optimization/milp\_solver.py}. The
Python code uses zero-based product indices internally, while product
identifiers in instance files may be one-based. The allowed route arcs contain
depot-to-station, station-to-station, and station-to-depot arcs, but not
depot-to-depot arcs. Empty depot-to-depot movements are priced only between
consecutive trips.

The formulation is exact for the selected trip bound \(R\). Increasing \(R\)
expands the feasible region, but it also increases the number of binary
variables and constraints. In computational experiments, \(R\) should therefore
be chosen as a compromise between operational flexibility and tractability.

\end{document}
